\subsection{Exact-commit hybrid and adaptive error composition}
\label{app:upper-proof-error}

This subsection completes the statistical part of the upper bound.
\Cref{app:lem:exact-hybrid} identifies a structurally safe exact-row hybrid
with the target law.  Separately,
\cref{app:lem:adaptive-hellinger-upper} bounds the distance between any
product-commit decision rule and its exact-row hybrid, whether or not its
batches are structurally safe.  A safe-decision-rule comparison then accounts
for failed screens in the implemented randomized sampler.


\paragraph{Structural safety.}
At a history $\History=(G,x_G)$, let $\ResidualForest$ be the true residual
forest on $\ResidualPositions=\PositionSet\setminus G$.
A complete decision rule is \emph{structurally safe} if, at every realized
history, each batch $B\subseteq\ResidualPositions$ it commits satisfies
\[
 |B\cap V(T)|\le1
 \qquad\text{for every connected component }T\text{ of }\ResidualForest.
\]
In particular, every singleton batch is safe.
This is a property used in the analysis, not permission for the implemented
sampler to inspect the hidden forest.

\paragraph{The two executions and their output laws.}
Fix one complete decision rule, the frozen oracle $\FrozenOracle$, and a
decision-rule seed $\ControllerSeed=w$, including the colors and all
auxiliary draft randomness.  The product-commit draws remain random after
this conditioning.  For a chosen batch $B$ and the history-compatible
commit state $y$ in \eqref{app:eq:commit-state}, compare:
\begin{itemize}
\item \emph{Oracle execution.}
A batch assignment $z\in\Vocab^B$ has mass
$\prod_{j\in B}\OracleRowMass{y}{j}{z_j}$.
Recall from \textcolor{red}{\cref{app:def:output-risk}} that its full output law is
$\SamplerOutputLaw{w}{\FrozenOracle}(x)
 =\Pr(\SamplerOutput=x\mid\ControllerSeed=w)$, with $\FrozenOracle$ fixed
and all product-commit draws integrated out.
\item \emph{Exact-commit hybrid.}
Replace only this commit mass by
$\prod_{j\in B}\ExactRowMass{y}{j}{z_j}$, using the exact rows in
\eqref{app:eq:exact-row}.  Keep the decision rule, draft generator, query
rules, and all frozen-oracle replies unchanged.
Denote this execution's full output law by $\ExactCommitOutputLaw{w}$.
\end{itemize}
Both are laws on $\Vocab^N$; $\otimes$ refers to the product within each
commit, not to independence of all output coordinates.
The hybrid can still depend on $\FrozenOracle$ through its decisions;
that dependence is suppressed in $\ExactCommitOutputLaw{w}$.
The shared rules do not imply shared realized transcripts: different commit
outcomes can lead to different later queries and batches.

Structural safety identifies the hybrid's product at each commit with the
target's joint batch conditional; for an unsafe batch they need not agree.
Thus \cref{app:lem:exact-hybrid} uses safety to prove
$\ExactCommitOutputLaw{w}=\TargetLaw$, whereas
\cref{app:lem:adaptive-hellinger-upper} compares the two output laws without
requiring safety.  Both lemmas also apply to the analysis-only decision rule
introduced below.

\paragraph{Intuition.}
A safe batch takes at most one vertex from each independent residual
component.  Its product of exact rows is therefore the exact joint
conditional.  Multiplying these joint conditionals along the adaptive
partition is the ordinary chain rule, so the output law is exactly
$\TargetLaw$.

\begin{lemma}[Exact hybrid equals the target law]
\label{app:lem:exact-hybrid}
For every structurally safe complete decision rule, the exact-commit hybrid
satisfies
\begin{equation}
 \ExactCommitOutputLaw{w}=\TargetLaw
 \qquad\text{for every decision-rule seed }w.
 \label{app:eq:exact-hybrid-target}
\end{equation}
\end{lemma}

\begin{proof}

\emph{1. Identify the exact joint law of a safe batch.}
For a singleton batch the product of exact singleton rows is trivially its
exact joint conditional.  For a nonsingleton safe batch, its vertices lie in
distinct true residual components.  Lemma~\ref{app:lem:residual-markov}
therefore identifies their exact joint conditional with the product of their
exact singleton conditionals, as in
\eqref{app:eq:centroid-batch-factorization}.


\emph{2. Apply the chain rule along each adaptive path.}
For fixed $w$, draft generation is deterministic as a function of the observed
past, even when its rule is randomized before conditioning on $w$.
For a full assignment $x\in\Vocab^N$, let $B_t(x)$ be the batch selected along
the unique hybrid path consistent with $x$, and let $G_{t-1}(x)$ be the set
committed before that batch, with $G_0(x)=\varnothing$ and
$G_t(x)=G_{t-1}(x)\cup B_t(x)$.  The batches form an adaptive ordered partition
of $\PositionSet$.  The preceding paragraph identifies every hybrid batch law
with the corresponding exact joint conditional, so the conditional chain
rule gives
\begin{align}
 \ExactCommitOutputLaw{w}(x)
 &=\prod_t
 \TargetLaw\!\left(
  X_{B_t(x)}=x_{B_t(x)}
 \mid X_{G_{t-1}(x)}=x_{G_{t-1}(x)}
 \right)\notag\\
 &=\prod_t
 \frac{\TargetLaw(X_{G_t(x)}=x_{G_t(x)})}
      {\TargetLaw(X_{G_{t-1}(x)}=x_{G_{t-1}(x)})}\notag\\
 &=\TargetLaw(x).
 \label{app:eq:adaptive-exact-chain-rule}
\end{align}
The batch partition may depend on earlier committed values, but for each fixed
$x$ the displayed factors are precisely the chain-rule factors along that
path.  Consecutive numerators and denominators cancel; the first denominator
is the probability of the empty constraint, hence one, and the last numerator
is $\TargetLaw(x)$ because every coordinate has been committed.
This proves the equality pointwise.
\end{proof}

We next compare the full output laws $\ExactCommitOutputLaw{w}$ (exact-row
commits) and $\SamplerOutputLaw{w}{\FrozenOracle}$ (oracle-row commits) for
the same decision rule and fixed seed.  For probability masses $P,Q$ on a
finite set, write
\begin{equation}
 \HellingerAffinity(P,Q)
 :=\sum_x\sqrt{P(x)Q(x)}
 =1-\sqHellinger(P,Q).
 \label{app:eq:affinity-definition}
\end{equation}

\paragraph{Intuition.}
The error budget is attached to committed coordinates, not to probes or
rounds.  Each coordinate is committed once, so the total budget is
\[
 N\cdot\frac{(\OracleRadius)^2}{2N}=\frac{(\OracleRadius)^2}{2}.
\]
For adaptive batches, backward induction on the execution tree justifies
this accounting.  At each node, charge the current batch and then use the
budget of the remaining coordinates.

\begin{lemma}[Adaptive squared-Hellinger composition]
\label{app:lem:adaptive-hellinger-upper}
{\color{red}Assume \cref{app:ass:oracle-accuracy} (A2): every valid row pair obeys
$\sqHellinger(\ExactRow{y}{j},\OracleRow{y}{j})
 \le(\OracleRadius)^2/(2N)$.  For every complete decision rule using the
product-commit rule and every fixed decision-rule seed $w$, regardless of
structural safety,}
\begin{equation}
 \sqHellinger\bigl(
  \ExactCommitOutputLaw{w},
  \SamplerOutputLaw{w}{\FrozenOracle}
 \bigr)
 \le\frac{(\OracleRadius)^2}{2}.
 \label{app:eq:global-upper-hellinger}
\end{equation}
\end{lemma}

\begin{proof}

\emph{1. Compare the two kernels at a common node.}
Fix $\FrozenOracle$ and $w$ throughout the proof.
\hypertarget{app-execution-tree-state}{}
An \emph{execution-tree node} is a stored algorithm state immediately before
a commit, or a terminal state after completion; it is not a vertex of the
hidden forest.  A pre-commit state records
\[
 \mathfrak h=(\History,\ScreenTranscript,\text{decision-rule state}),
 \qquad
 \History=(G,x_G),\quad
 \ResidualPositions=\PositionSet\setminus G.
\]
Here $\ScreenTranscript$ contains the preceding query/reply and commit
record, and the internal state includes any current draft and counters.
We use this complete stored state to define continuation kernels directly,
without replaying the decision rule from its initial state.
With $\FrozenOracle$ and $w$ fixed, all operations between commits are
deterministic.  Thus, from the \emph{same stored state} $\mathfrak h$, both
executions receive the same frozen-oracle replies and select the same
batch $B(\mathfrak h)\subseteq\ResidualPositions$ and commit state
$y(\mathfrak h)$ from \eqref{app:eq:commit-state}.
Comparing kernels at this common state does not assert that two separately
sampled executions have identical paths.
For a batch assignment $z\in\Vocab^{B(\mathfrak h)}$, their kernels are
\begin{align}
 \ExactBatchMass{\mathfrak h}{z}
 &:=\prod_{j\in B(\mathfrak h)}
       \ExactRowMass{y(\mathfrak h)}{j}{z_j},\notag\\
 \OracleBatchMass{\mathfrak h}{z}
 &:=\prod_{j\in B(\mathfrak h)}
       \OracleRowMass{y(\mathfrak h)}{j}{z_j}.
 \label{app:eq:exact-oracle-batch-laws}
\end{align}
Both kernels are probability masses on $\Vocab^{B(\mathfrak h)}$:
\begin{align*}
 \sum_{z\in\Vocab^{B(\mathfrak h)}}\ExactBatchMass{\mathfrak h}{z}
 &=\prod_{j\in B(\mathfrak h)}
       \sum_{a\in\Vocab}\ExactRowMass{y(\mathfrak h)}{j}{a}=1,\\
 \sum_{z\in\Vocab^{B(\mathfrak h)}}\OracleBatchMass{\mathfrak h}{z}
 &=\prod_{j\in B(\mathfrak h)}
       \sum_{a\in\Vocab}\OracleRowMass{y(\mathfrak h)}{j}{a}=1.
\end{align*}
The finite sums factor by distributivity over the Cartesian product.
The first kernel is a product of exact rows, whether or not it equals the
target's joint batch conditional.  \textcolor{red}{Hence this comparison needs
no structural safety condition.}


\emph{2. Charge the current batch to its coordinates.}
Apply the product identity in \cref{app:lem:finite-hellinger-calculus}
with index set $B(\mathfrak h)$, each coordinate space equal to $\Vocab$,
and component laws $\ExactRow{y(\mathfrak h)}{j}$ and
$\OracleRow{y(\mathfrak h)}{j}$.  This gives
\begin{equation}
 \HellingerAffinity(\ExactBatchLaw{\mathfrak h},
                    \OracleBatchLaw{\mathfrak h})
 =\prod_{j\in B(\mathfrak h)}
  \left[1-\sqHellinger\bigl(
        \ExactRow{y(\mathfrak h)}{j},
        \OracleRow{y(\mathfrak h)}{j}\bigr)\right].
 \label{app:eq:upper-batch-affinity}
\end{equation}
Using $1-\prod_j(1-u_j)\le\sum_j u_j$ for $u_j\in[0,1]$ and
\textcolor{red}{\eqref{app:eq:A2}},
\begin{align}
 1-\HellingerAffinity(\ExactBatchLaw{\mathfrak h},
                      \OracleBatchLaw{\mathfrak h})
 &=\sqHellinger(\ExactBatchLaw{\mathfrak h},
                \OracleBatchLaw{\mathfrak h})
 \label{app:eq:local-batch-cost-definition}\\
 &\le\sum_{j\in B(\mathfrak h)}
     \sqHellinger\bigl(\ExactRow{y(\mathfrak h)}{j},
                      \OracleRow{y(\mathfrak h)}{j}\bigr)\notag\\
 &\le |B(\mathfrak h)|\frac{(\OracleRadius)^2}{2N}.
 \label{app:eq:local-batch-cost-bound}
\end{align}


\emph{3. Derive the suffix-affinity recursion.}
Take a stored execution state $\mathfrak h$ from part~1 of this proof, with committed
history $\History=(G,x_G)$ and remaining set
$\ResidualPositions=\PositionSet\setminus G$.
Here a \emph{suffix} is the assignment of all still-uncommitted coordinates,
not necessarily a contiguous block of positions.  Define:
\begin{itemize}
\item $\ExactSuffixLaw{\mathfrak h}$: the law of those remaining values
when execution resumes from this entire stored state using exact-row commits.
\item $\OracleSuffixLaw{\mathfrak h}$: the corresponding law when it
resumes from the same state using oracle-row commits.
\end{itemize}
Both are probability laws on $\Vocab^{\ResidualPositions}$, with
$\FrozenOracle$ and $w$ still fixed.  The recursion below defines these
continuation laws from the batch kernels in
\eqref{app:eq:exact-oracle-batch-laws}; at a terminal state
$\ResidualPositions=\varnothing$, each puts mass one on the empty assignment.
Set
\begin{equation}
 \SuffixAffinity{\mathfrak h}
 :=\HellingerAffinity(\ExactSuffixLaw{\mathfrak h},
                      \OracleSuffixLaw{\mathfrak h}).
 \label{app:eq:suffix-affinity}
\end{equation}
After a batch value $z\in\Vocab^{B(\mathfrak h)}$, write
$\mathfrak h\oplus z$ for the next pre-commit node (or terminal node),
including all intervening deterministic operations.  It has remaining
coordinates $\ResidualPositions\setminus B(\mathfrak h)$.
For $x'\in\Vocab^{\ResidualPositions\setminus B(\mathfrak h)}$,
the suffix assignment $(z,x')$ has masses
\begin{align*}
 \ExactSuffixMass{\mathfrak h}{(z,x')}
 &=\ExactBatchMass{\mathfrak h}{z}\,
   \ExactSuffixMass{\mathfrak h\oplus z}{x'},\\
 \OracleSuffixMass{\mathfrak h}{(z,x')}
 &=\OracleBatchMass{\mathfrak h}{z}\,
   \OracleSuffixMass{\mathfrak h\oplus z}{x'}.
\end{align*}
These are normalized: summing first over $x'$ gives one for each child
law, and summing over $z$ then gives one by the batch normalization above.
There are at most $N$ nonempty commits, so backward induction starts at
the empty suffix and defines every law.
Summing the geometric means first over $x'$ now yields
\begin{align}
 \SuffixAffinity{\mathfrak h}
 &=\sum_z\sqrt{\ExactBatchMass{\mathfrak h}{z}
               \OracleBatchMass{\mathfrak h}{z}}
       \sum_{x'}\sqrt{
        \ExactSuffixMass{\mathfrak h\oplus z}{x'}
        \OracleSuffixMass{\mathfrak h\oplus z}{x'}}\notag\\
 &=\sum_z\sqrt{\ExactBatchMass{\mathfrak h}{z}
               \OracleBatchMass{\mathfrak h}{z}}
       \SuffixAffinity{\mathfrak h\oplus z}.
 \label{app:eq:suffix-affinity-recursion}
\end{align}
At a terminal node the suffix is empty and its affinity is one.


\emph{4. Bound the remaining loss by backward induction.}
At a node with committed set $G$, every complete continuation partitions
$\ResidualPositions$ into disjoint future batches.  Thus its total
coordinate budget is
\begin{equation}
 \sum_{\text{future }t}|B(\mathfrak h_t)|
       \frac{(\OracleRadius)^2}{2N}
 =|\ResidualPositions|\frac{(\OracleRadius)^2}{2N}.
 \label{app:eq:remaining-path-budget}
\end{equation}
We prove by backward induction that this budget bounds the suffix loss:
\begin{equation}
 1-\SuffixAffinity{\mathfrak h}
 \le|\ResidualPositions|\frac{(\OracleRadius)^2}{2N}.
 \label{app:eq:adaptive-affinity-induction}
\end{equation}
The terminal case has both sides zero.  At a nonterminal node, subtract
the recursion from one and add and subtract the batch affinity:
\begin{align*}
 1-\SuffixAffinity{\mathfrak h}
 &=1-\HellingerAffinity(\ExactBatchLaw{\mathfrak h},
                        \OracleBatchLaw{\mathfrak h})\\
 &\quad+\sum_z
   \sqrt{\ExactBatchMass{\mathfrak h}{z}\OracleBatchMass{\mathfrak h}{z}}
        \bigl(1-\SuffixAffinity{\mathfrak h\oplus z}\bigr)\\
 &\le |B(\mathfrak h)|\frac{(\OracleRadius)^2}{2N}
   +\bigl(|\ResidualPositions|-|B(\mathfrak h)|\bigr)
      \frac{(\OracleRadius)^2}{2N}
      \sum_z\sqrt{\ExactBatchMass{\mathfrak h}{z}
                  \OracleBatchMass{\mathfrak h}{z}}\\
 &\le|\ResidualPositions|\frac{(\OracleRadius)^2}{2N}.
\end{align*}
The first inequality uses the batch bound and the induction hypothesis:
every child has exactly $|\ResidualPositions|-|B(\mathfrak h)|$
uncommitted coordinates.  \textcolor{red}{In the last inequality, the
remaining sum is the batch affinity, at most $\sqrt{1\cdot1}=1$ by
Cauchy--Schwarz and the two kernel normalizations in part~1.  The resulting
coordinate count is $|B(\mathfrak h)|+(|\ResidualPositions|-|B(\mathfrak h)|)
=|\ResidualPositions|$.}


\emph{5. Evaluate the budget at the root.}
The root $\mathfrak h_{\emptyset}$ is the common state before the first
commit, after the initial deterministic operations.  It has $G=\varnothing$
and $\ResidualPositions=\PositionSet$, so its continuation laws are the
full output laws recalled above:
\[
 \ExactSuffixLaw{\mathfrak h_{\emptyset}}=\ExactCommitOutputLaw{w},
 \qquad
 \OracleSuffixLaw{\mathfrak h_{\emptyset}}
 =\SamplerOutputLaw{w}{\FrozenOracle}.
\]
In particular, $|\ResidualPositions|=N$.  Equivalently, along every complete
path the disjoint commit batches satisfy
\begin{equation}
 \sum_t|B(\mathfrak h_t)|=N.
 \label{app:eq:path-batches-partition-N}
\end{equation}
Therefore
\begin{equation}
 \sqHellinger(\ExactCommitOutputLaw{w},
              \SamplerOutputLaw{w}{\FrozenOracle})
 =1-\SuffixAffinity{\mathfrak h_{\emptyset}}
 \le N\frac{(\OracleRadius)^2}{2N}
 =\frac{(\OracleRadius)^2}{2}.
 \label{app:eq:root-path-budget}
\end{equation}
Counterfactual operations affect which common node is reached; they do not
generate committed coordinates and add no term to this budget.
\end{proof}

We now separate oracle error from screen failures by comparing with an
analysis-only safe decision rule, whose seed-averaged discrepancy from the
sampler is charged to the failed-screen probability.

Recall the \hyperlink{app-reached-screen-failure}{reached-call event}
$E_s=\{R_s\text{ occurs and }\ScreenSuccess{H^{(s)}}\text{ fails}\}$:
$R_s$ means the implemented sampler reaches call $s$, and $H^{(s)}$ is
its history.  Failure means that the separating-color majority condition
\eqref{app:eq:screen-success-event} is violated.
Thus $\ScreenFailureEvent=\bigcup_{s=1}^{\ScreenCallCap}E_s$.

\paragraph{Intuition.}
The safe decision rule follows the sampler until a screen fails and then
finishes by singleton commits.  Its batches are always safe, so the previous
two lemmas bound its TV error by $\OracleRadius$.  A common-path coupling
charges the discrepancy between the two outputs only to screen failure:
\[
 \mathbb E_W\mathrm{TV}(\hbox{target},\hbox{sampler})
 \le\OracleRadius+\ScreenFailureBudget.
\]
The first contribution is controlled for every fixed seed; the failure
contribution is bounded only after averaging over the colors and commit
draws, as verified below.
\textcolor{red}{We compare the full output laws: conditioning the implemented
output on screen success would also condition its reached histories.}

\begin{lemma}[Safe decision rule and failed-screen coupling]
\label{app:lem:safe-controller-comparison}
Fix a target--oracle pair satisfying the finite hypotheses of
Theorem~\ref{app:thm:finite-upper}.  There is an analysis-only structurally
safe decision rule whose conditional output law
$\SafeSamplerOutputLaw{w}{\FrozenOracle}$ satisfies
\begin{align}
 \dTV\bigl(\TargetLaw,
       \SafeSamplerOutputLaw{w}{\FrozenOracle}\bigr)
 &\le \OracleRadius
 &&\text{for every }w,
 \label{app:eq:safe-seedwise-TV}\\
 \mathbb E_{\ControllerSeed}
 \dTV\bigl(
   \SamplerOutputLaw{\ControllerSeed}{\FrozenOracle},
   \SafeSamplerOutputLaw{\ControllerSeed}{\FrozenOracle}
 \bigr)
 &\le\ScreenFailureBudget.
 \label{app:eq:implemented-safe-TV}
\end{align}
\end{lemma}

\begin{proof}

\emph{1. Define the analysis-only safe decision rule.}
The safe decision rule simulates Algorithm~\ref{app:alg:complete-sampler} with
the same fixed public row selector, draft generator, auxiliary seed, and fresh colorings while every
reached screen satisfies
\eqref{app:eq:screen-success-event}.  It is additionally allowed to inspect
the hidden residual forest.  Immediately after the first unsuccessful
execution of \AlgoBlockRef{screen}, and also before any \AlgoGuardRef\
fallback, it stops using the screen
output and commits all remaining coordinates as singletons in increasing
position order.  More formally, number the successive simulation checkpoints
(screen returns and pre-fallback checks) by $n$, and define
\[
 \tau_{\mathrm{sw}}
 :=\inf\left\{n:
 \begin{array}{l}
 \text{an unsuccessful screen has just returned, or}\\
 \text{the simulated sampler is about to activate a fallback}
 \end{array}\right\},
 \qquad \inf\varnothing:=\infty.
\]
Before this checkpoint the two rules agree; from this checkpoint onward
the safe rule uses only the singleton completion just specified.
The checkpoint index is not the peel-phase index.
Hidden-forest inspection defines only this comparison law; it is not
available to the implemented sampler.
Each singleton commit uses the same frozen oracle $\FrozenOracle$ at the
updated committed history.  Writing $\widehat X^{\mathrm{safe}}$ for its
output, define
\[
 \SafeSamplerOutputLaw{w}{\FrozenOracle}(x)
 :=\Pr(\widehat X^{\mathrm{safe}}=x\mid\ControllerSeed=w),
\]
with the target and oracle fixed and all product-commit draws integrated out.


\emph{2. Check structural safety in all three stages.}
\begin{itemize}
\item Before a switch, the peel commits in \AlgoBlockRef{peel} are singletons.
\item If the terminal stage is
reached without a switch, \cref{app:lem:successful-path-structure} identifies
the graph constructed in \AlgoBlockRef{terminal} with the true residual
forest, so every centroid batch in \AlgoBlockRef{centroid} contains one
vertex from each true residual component.
\item After a switch, all
batches are singletons.
\end{itemize}
Hence the safe decision rule is structurally safe on
every path.

\emph{3. Bound oracle error for each fixed seed.}
Apply \cref{app:lem:exact-hybrid,app:lem:adaptive-hellinger-upper}
to this safe rule and its own product-of-exact-rows hybrid.
The first lemma identifies that hybrid with $\TargetLaw$; the second
compares it with the safe rule's oracle output.  Thus, for each fixed $w$,
\begin{equation}
 \sqHellinger\bigl(
  \TargetLaw,\SafeSamplerOutputLaw{w}{\FrozenOracle}
 \bigr)
 \le\frac{(\OracleRadius)^2}{2}.
 \label{app:eq:safe-global-hellinger}
\end{equation}
By \eqref{app:eq:aux-tv-hellinger},
\begin{equation}
 \dTV\bigl(\TargetLaw,
       \SafeSamplerOutputLaw{w}{\FrozenOracle}\bigr)
 \le\sqrt{2\sqHellinger\bigl(
  \TargetLaw,\SafeSamplerOutputLaw{w}{\FrozenOracle}
 \bigr)}
 \le\OracleRadius,
\end{equation}
which proves \eqref{app:eq:safe-seedwise-TV}.


\emph{4. Couple the sampler to its own safe decision rule.}
For \eqref{app:eq:implemented-safe-TV}, fix the full decision-rule seed $w$ and couple the
two conditional executions as follows.  Before the safe decision rule switches,
whenever their histories agree, use
the same draw from their common product-commit law.  Thus their histories and
all drafts and deterministic frozen-oracle replies remain identical until the safe
decision rule switches.  By \cref{app:lem:upper-resources}, in the
implemented execution,
\[
 \{\text{cycle repair or a hard round-cap fallback occurs}\}
 \subseteq\ScreenFailureEvent.
\]
An empty peel set or the peel-phase cap triggers the normal loop exit in
\AlgoBlockRef{peel}.  On $\ScreenFailureEvent^{\mathsf c}$ no switch occurs
and the outputs coincide; repair can start only after the safe rule has
switched to singleton completion.  For the coupled outputs
$\SamplerOutput$ and $\widehat X^{\mathrm{safe}}$, the common prefix gives
the pathwise inclusions
\[
 \{\SamplerOutput\ne\widehat X^{\mathrm{safe}}\}
 \subseteq\{\tau_{\mathrm{sw}}<\infty\}
 \subseteq\ScreenFailureEvent.
\]
Indeed, a first switch caused by a failed screen occurs in that common
prefix.  A first switch caused by a fallback with no preceding failure is
excluded by \cref{app:lem:upper-resources}.  After a switch the two
executions may be completed with any coupling of their respective laws.
The first common-path failure is exactly the first failure of the
implemented execution, since the coupled histories agree up to that point.
\textcolor{red}{For any coupling $(X,Y)$ of laws $P,Q$, the coupling inequality
is $\dTV(P,Q)\le\Pr(X\ne Y)$.  Applying it to these two outputs gives}
{\color{red}
\begin{equation}
 \dTV\bigl(
   \SamplerOutputLaw{w}{\FrozenOracle},
   \SafeSamplerOutputLaw{w}{\FrozenOracle}
 \bigr)
 \le
 \Pr(\SamplerOutput\ne\widehat X^{\mathrm{safe}}
       \mid\ControllerSeed=w)
 \le
 \Pr(\ScreenFailureEvent\mid\ControllerSeed=w),
 \label{app:eq:conditional-safe-coupling}
\end{equation}
}
where the probability integrates the coupled commit draws.

\emph{5. Average the failed-screen probability over seeds.}
Averaging \eqref{app:eq:conditional-safe-coupling} over the seed gives
\begin{align*}
 &\mathbb E_W\dTV\bigl(\SamplerOutputLaw{W}{\FrozenOracle},
                        \SafeSamplerOutputLaw{W}{\FrozenOracle}\bigr)\\
 &\quad\le\mathbb E_W\Pr(\ScreenFailureEvent\mid W)
 =\Pr(\ScreenFailureEvent)
 =\Pr\!\left(\bigcup_{s=1}^{\ScreenCallCap}E_s\right)
 \le\ScreenFailureBudget.
\end{align*}
The equality averaging conditional probabilities is the law of total
expectation over colors, draft randomness, and commit draws; the final
inequality is \eqref{app:eq:adaptive-screen-failure}.  This proves
\eqref{app:eq:implemented-safe-TV}, without a seedwise failure bound.
\end{proof}

\begin{proof}[Proof of Theorem~\ref{app:thm:finite-upper}]

\emph{1. Account for submissions, commits, and depth.}
With the standard zero-query generator, the pathwise query and round bounds are
\eqref{app:eq:proof-finite-Q}--\eqref{app:eq:proof-finite-R} from
\cref{app:lem:upper-resources}.  The all-mask preprocessing state is excluded
from $\CounterfactualQueries$ by
\textcolor{red}{\cref{app:def:resources}}. Charging it gives
$\ReportedQueries=\CounterfactualQueries+1$, as in
\eqref{app:eq:reported-query-count}.
The depth bound is \eqref{app:eq:finite-oracle-depth}.


\emph{2. Combine oracle error and failed-screen error.}
For every target--oracle pair, the TV triangle inequality followed by
\cref{app:lem:safe-controller-comparison} gives
\begin{align}
 &\mathbb E_{\ControllerSeed}
 \dTV\bigl(
  \TargetLaw,
  \SamplerOutputLaw{\ControllerSeed}{\FrozenOracle}
 \bigr)\notag\\
 &\quad\le
 \mathbb E_{\ControllerSeed}
 \dTV\bigl(
  \TargetLaw,
  \SafeSamplerOutputLaw{\ControllerSeed}{\FrozenOracle}
 \bigr)
 +\mathbb E_{\ControllerSeed}
 \dTV\bigl(
  \SafeSamplerOutputLaw{\ControllerSeed}{\FrozenOracle},
  \SamplerOutputLaw{\ControllerSeed}{\FrozenOracle}
 \bigr)\notag\\
 &\quad\le
 \OracleRadius+\ScreenFailureBudget.
 \label{app:eq:randomized-upper-TV-proof}
\end{align}
This proves \eqref{app:eq:finite-upper-seed-risk} for each fixed oracle.
Taking the supremum over
$\FrozenOracle\in\OracleClass(\TargetLaw;\OracleRadius)$ gives
\eqref{app:eq:finite-upper-risk}.  With the choice
\eqref{app:eq:finite-failure-budget-choice}, its
right-hand side is $3\OracleRadius/2$.
\end{proof}

% Please keep the selector-invariance supplement in
% src/sections/proofs/row_selector_invariance.tex, included by the auxiliary appendix.
